\thispagestyle{plain}
\phantomsection
\addcontentsline{toc}{chapter}{中文摘要}
\centerline{\zihao{-3}\heiti 中文摘要}

\linespread{1.4}\zihao{-4}\bigskip

本文针对自动化车床刀具故障场景，研究“检查间隔 + 换刀策略”的联合优化问题。依据题目给出的 100 组刀具寿命样本与费用参数，建立以单位合格品平均成本最小为目标的期望成本模型。模型同时考虑废品损失、检查成本、故障恢复成本、主动换刀成本，并在第二问中纳入误判停机损失。

在方法上，先利用刀具寿命数据估计寿命分布参数，再将检查策略参数化为检查序列，构造分层优化：固定检查次数求最优检查点，再在不同检查次数之间选取全局最优。针对两种质量判定情形分别求解，并比较策略差异。进一步讨论双检确认与后段加密检查等改进方案，以降低误判带来的额外损失。

研究表明，该问题可统一为约束非线性优化框架；最优策略与故障判定准确率、检查成本和换刀成本显著相关。该建模过程形成了“数据分布拟合 -- 期望损失建模 -- 策略优化 -- 灵敏度分析”的完整流程，可为离散制造系统中的预防性维护与质量控制提供参考。

\bigskip
\noindent{\zihao{4}\heiti 关键词：}
刀具寿命；检查策略；换刀策略；期望成本；非线性优化
